
\documentclass{article}
\usepackage{colortbl}
\usepackage{makecell}
\usepackage{multirow}
\usepackage{supertabular}

\begin{document}

\newcounter{utterance}

\centering \large Interaction Transcript for game `wordle', experiment `medium\_frequency\_words\_no\_clue\_no\_critic', episode 2 with refineonly{-}Qwen3.5{-}4B.
\vspace{24pt}

{ \footnotesize  \setcounter{utterance}{1}
\setlength{\tabcolsep}{0pt}
\begin{supertabular}{c@{$\;$}|p{.15\linewidth}@{}p{.15\linewidth}p{.15\linewidth}p{.15\linewidth}p{.15\linewidth}p{.15\linewidth}}
    \# & \multicolumn{2}{c}{Player} && \multicolumn{2}{c}{Game Master} \\
    \hline

    \theutterance \stepcounter{utterance}  
    & & & \multicolumn{4}{p{0.6\linewidth}}{
        \cellcolor[rgb]{0.9,0.9,0.9}{
            \makecell[{{p{\linewidth}}}]{
                \texttt{\tiny{[P1$\langle$GM]}}
                \texttt{You are a language wizard who likes to guess words by using the given rules.} \\
\\ 
\texttt{Welcome to Wordle! You have six attempts to guess the target word, a valid English word of five lowercase letters (a{-}z). Please use the tags "explanation:" and "guess:" to provide a concise explanation for each guess.} \\
\\ 
\texttt{For instance, if your guess is "apple", your response should be} \\
\texttt{explanation: this is a common five{-}letter English word, and I am starting my guess with this word.} \\
\texttt{guess: apple} \\
\\ 
\texttt{After each guess, your answer will be validated, and you will receive feedback indicating which letters are correct (green), which letters are correct but in the wrong position (yellow), and which letters are incorrect (red). This feedback can be useful in determining which letters to include or exclude in your next guess.} \\
\\ 
\texttt{For example, the feedback for "apple" might be:} \\
\texttt{guess\_feedback: a<yellow> p<yellow> p<green> l<yellow> e<red>} \\
\\ 
\texttt{The explanation should contain details about how the guess\_feedback is used to arrive at a new guess.} \\
\\ 
\texttt{Do not generate any other text.} \\
\\ 
\texttt{Let's begin with your first guess.} \\
            }
        }
    }
     \\ \\

    \theutterance \stepcounter{utterance}  
    & \multicolumn{4}{p{0.6\linewidth}}{
        \cellcolor[rgb]{0.561,0.737,0.561}{
            \makecell[{{p{\linewidth}}}]{
                \texttt{\tiny{[P1$\rangle$GM]}}
                \texttt{explanation: I will start with a very common five{-}letter word that contains a variety of frequent letters to maximize information gain.} \\
\texttt{guess: apple} \\
            }
        }
    }
    & & \\ \\

    \theutterance \stepcounter{utterance}  
    & & & \multicolumn{2}{p{0.3\linewidth}}{
        \cellcolor[rgb]{0.9,0.9,0.9}{
            \makecell[{{p{\linewidth}}}]{
                \texttt{\tiny{[GM$|$GM]}}
                \texttt{attempts = 1} \\
\texttt{target\_word = resin} \\
\texttt{guess = apple} \\
\texttt{guess\_feedback = a<red> p<red> p<red> l<red> e<yellow>} \\
            }
        }
    }
    & & \\ \\

    \theutterance \stepcounter{utterance}  
    & & & \multicolumn{4}{p{0.6\linewidth}}{
        \cellcolor[rgb]{0.9,0.9,0.9}{
            \makecell[{{p{\linewidth}}}]{
                \texttt{\tiny{[P1$\langle$GM]}}
                \texttt{guess\_feedback: a<red> p<red> p<red> l<red> e<yellow>} \\
\\ 
\texttt{Provide your response only in this format.} \\
\texttt{explanation: details} \\
\texttt{guess: word} \\
            }
        }
    }
     \\ \\

    \theutterance \stepcounter{utterance}  
    & \multicolumn{4}{p{0.6\linewidth}}{
        \cellcolor[rgb]{0.561,0.737,0.561}{
            \makecell[{{p{\linewidth}}}]{
                \texttt{\tiny{[P1$\rangle$GM]}}
                \texttt{explanation: The feedback indicates that 'e' is in the word but in the last position, while 'a', 'p', and 'l' are not present. I need to find a word ending in 'e' that does not contain 'a', 'p', or 'l'. I will try a common word ending in 'e' with different vowels and consonants.} \\
\texttt{guess: crane} \\
            }
        }
    }
    & & \\ \\

    \theutterance \stepcounter{utterance}  
    & & & \multicolumn{2}{p{0.3\linewidth}}{
        \cellcolor[rgb]{0.9,0.9,0.9}{
            \makecell[{{p{\linewidth}}}]{
                \texttt{\tiny{[GM$|$GM]}}
                \texttt{attempts = 2} \\
\texttt{target\_word = resin} \\
\texttt{guess = crane} \\
\texttt{guess\_feedback = c<red> r<yellow> a<red> n<yellow> e<yellow>} \\
            }
        }
    }
    & & \\ \\

    \theutterance \stepcounter{utterance}  
    & & & \multicolumn{4}{p{0.6\linewidth}}{
        \cellcolor[rgb]{0.9,0.9,0.9}{
            \makecell[{{p{\linewidth}}}]{
                \texttt{\tiny{[P1$\langle$GM]}}
                \texttt{guess\_feedback: c<red> r<yellow> a<red> n<yellow> e<yellow>} \\
\\ 
\texttt{Provide your response only in this format.} \\
\texttt{explanation: details} \\
\texttt{guess: word} \\
            }
        }
    }
     \\ \\

    \theutterance \stepcounter{utterance}  
    & \multicolumn{4}{p{0.6\linewidth}}{
        \cellcolor[rgb]{0.561,0.737,0.561}{
            \makecell[{{p{\linewidth}}}]{
                \texttt{\tiny{[P1$\rangle$GM]}}
                \texttt{explanation: The feedback from "crane" (c<red> r<yellow> a<red> n<yellow> e<yellow>) indicates that 'r', 'n', and 'e' are in the word, but not in the 2nd, 4th, and 5th positions respectively. Since 'a' is red, I must avoid it. I need a 5{-}letter word containing r, n, e where 'e' is not at the end, 'r' is not at position 2, and 'n' is not at position 4. "sneer" fits well: it contains s, n, e, e, r. In "sneer", 'n' is at position 2 (which contradicts the yellow 'n' at position 4 in "crane" only if the target has 'n' elsewhere, but wait: in "crane", 'n' was at pos 4 and got yellow, meaning 'n' is in the word but not at pos 4. In "sneer", 'n' is at pos 2. If the target is "sneer", then 'n' is at pos 2. In "crane", 'n' was at pos 4. Since 'n' is yellow in "crane", it means 'n' is in the word but not at pos 4. This is consistent with "sneer" where 'n' is at pos 2. Also 'r' in "crane" was at pos 2 (yellow), so 'r' is not at pos 2. In "sneer", 'r' is at pos 5. This is consistent. 'e' in "crane" was at pos 5 (yellow), so 'e' is not at pos 5. In "sneer", 'e' is at pos 2 and 4. This is consistent. So "sneer" is a strong candidate.} \\
\texttt{guess: sneer} \\
            }
        }
    }
    & & \\ \\

    \theutterance \stepcounter{utterance}  
    & & & \multicolumn{2}{p{0.3\linewidth}}{
        \cellcolor[rgb]{0.9,0.9,0.9}{
            \makecell[{{p{\linewidth}}}]{
                \texttt{\tiny{[GM$|$GM]}}
                \texttt{attempts = 3} \\
\texttt{target\_word = resin} \\
\texttt{guess = sneer} \\
\texttt{guess\_feedback = s<yellow> n<yellow> e<yellow> e<red> r<yellow>} \\
            }
        }
    }
    & & \\ \\

    \theutterance \stepcounter{utterance}  
    & & & \multicolumn{4}{p{0.6\linewidth}}{
        \cellcolor[rgb]{0.9,0.9,0.9}{
            \makecell[{{p{\linewidth}}}]{
                \texttt{\tiny{[P1$\langle$GM]}}
                \texttt{guess\_feedback: s<yellow> n<yellow> e<yellow> e<red> r<yellow>} \\
\\ 
\texttt{Provide your response only in this format.} \\
\texttt{explanation: details} \\
\texttt{guess: word} \\
            }
        }
    }
     \\ \\

    \theutterance \stepcounter{utterance}  
    & \multicolumn{4}{p{0.6\linewidth}}{
        \cellcolor[rgb]{0.561,0.737,0.561}{
            \makecell[{{p{\linewidth}}}]{
                \texttt{\tiny{[P1$\rangle$GM]}}
                \texttt{explanation: The feedback for "sneer" (s<yellow> n<yellow> e<yellow> e<red> r<yellow>) tells us the letters are s, n, e, r. The 'e' at position 4 is red, meaning the 4th letter is not 'e'. The 'e' at position 3 is yellow, so the 3rd letter is 'e'. The 'r' at position 5 is yellow, so the 5th letter is not 'r'. The 'n' at position 2 is yellow, so the 2nd letter is not 'n'. The 's' at position 1 is yellow, so the 1st letter is not 's'. We have the pattern \_ \_ e \_ \_. We know 'r' must be in the word but not at pos 5. 'n' must be in the word but not at pos 2. 's' must be in the word but not at pos 1. Since we have 4 letters (s, n, e, r) and 5 slots, one letter is missing or one is repeated. Given previous feedback, we likely have one 'e'. We need to place s, n, r into slots 1, 2, 4, 5 such that s!=1, n!=2, r!=5, and slot 3 is 'e'. Let's try placing 'r' at position 4: \_ \_ e r \_. Now we need to place 's' and 'n' in slots 1 and 2. Since s cannot be 1 and n cannot be 2, this seems impossible if we only have one of each. Wait, maybe 's' is at 2 and 'n' is at 1? That gives "nsere". Is that a word? No. Maybe I missed a letter or the count is different. Let's reconsider the letters. Previous guesses: apple (e yellow), crane (r yellow, n yellow, e yellow), sneer (s yellow, n yellow, e yellow, r yellow). Letters found: s, n, e, r. Positions known: 3rd is 'e'. 2nd is not 'n'. 1st is not 's'. 5th is not 'r'. 4th is not 'e'. We need a 5th letter. Maybe 'o'? "snore" had 'o' but 'a' was red in crane. Wait, in "crane", 'a' was red. In "sneer", 'a' is not present. So 'a' is out. What about "sneer" itself? No, 4th letter was red. What about "sneak"? 'k' never seen. "snare"? 'a' is red. "snore"? 'o} \\
            }
        }
    }
    & & \\ \\

    \theutterance \stepcounter{utterance}  
    & & & \multicolumn{2}{p{0.3\linewidth}}{
        \cellcolor[rgb]{0.9,0.9,0.9}{
            \makecell[{{p{\linewidth}}}]{
                \texttt{\tiny{[GM$|$GM]}}
                \texttt{Error: The response should contain the 'guess:' keyword exactly once.} \\
            }
        }
    }
    & & \\ \\

    \theutterance \stepcounter{utterance}  
    & & & \multicolumn{2}{p{0.3\linewidth}}{
        \cellcolor[rgb]{0.9,0.9,0.9}{
            \makecell[{{p{\linewidth}}}]{
                \texttt{\tiny{[GM$|$GM]}}
                \texttt{game\_result = ABORT} \\
            }
        }
    }
    & & \\ \\

\end{supertabular}
}

\end{document}
